\chapter{Discussion and Future work}

\startcontents[chapters]

Attention, a concept extensively explored in cognitive science and machine learning, has made significant strides in the fields of computer vision and natural language processing. Self-attention-based architectures, prevalent in these domains, have demonstrated remarkable performance on various benchmarks. The self-attention mechanism employed by Transformer networks bears resemblance to the central executive controller attention theory proposed in cognitive psychology, as both are concerned with attentional control and cognitive processing. According to the central executive controller attention theory, attentional control relies on a central executive system that monitors and governs the flow of information within working memory. This system selectively attends to different components of the working memory representation based on task demands and information relevance.

Similarly, the self-attention mechanism in Transformer networks enables models to selectively attend to specific elements of the input sequence, guided by their relevance to the current task. This mechanism can be viewed as an abstraction of the central executive system, allowing the model to dynamically adjust its attentional focus in response to the input and its current state.

Furthermore, both the central executive system and the self-attention mechanism in Transformers involve the integration of multiple information sources, including internal representations and external cues, to facilitate attentional control and cognitive processing. In Transformers, this integration is achieved through multiple self-attention layers, enabling the model to construct an internal representation of the input sequence that captures the most pertinent features for the task at hand. Although the self-attention mechanism in Transformers and the central executive system in cognitive psychology are not identical, they share similar principles of attentional control and cognitive processing, enabling flexible and adaptive behavior to address evolving task demands and environmental cues.

Contrasting attention in Transformers with attention in the human mind reveals several distinguishing characteristics. In Transformers, attention "heads" operate independently in relation to the previous layer, indicating a lack of integration and coherence when compared to the integrated attention observed in the human mind. Transformers exhibit offline processing during learning, focusing solely on the present configuration, while attention in the human mind is influenced by past states, shaping the interpretation of the present and future.

Additionally, representations in Transformers are static, whereas in the human mind, representations are dynamic. Transformers employ features for scene segmentation, while in the human mind, features serve as parametric operators that facilitate scene predictions. Moreover, attention in Transformers does not drive active perception, unlike attention in the human mind, which plays a crucial role in guiding active perception and goal-directed cognition. Besides, attention in the human mind integrates latent features into a scene graph, enhancing the coherence and organization of information, a characteristic not explicitly addressed in the context of Transformers. These comparisons highlight the distinctions between attention mechanisms in Transformers and the human mind, emphasizing the integrated, dynamic, and goal-directed nature of attention within the human cognitive system.

Furthermore, attention, as a cognitive process enabling focused concentration on relevant stimuli, plays a vital role in enhancing human reasoning ability. To deepen our understanding of the self-attention mechanism, this thesis investigates its role in cognitive and computer vision architectures, particularly within the domain of visual reasoning. By exploring the interaction between self-attention and cognitive processes, we aim to contribute to the advancement of both cognitive science and machine learning, ultimately enhancing our comprehension and utilization of attention mechanisms.

Visual reasoning is the process of analyzing the provided visual information in order to solve a task. It is considered an important part of fluid intelligence, which involves thinking and reasoning independent of learning, education, and experience. This ability has not only been shown in primates \citep{gentner2021learning} but also in bees \citep{giurfa2001concepts} and in newborn ducklings \citep{martinho2016ducklings}. On the contrary, prior studies \citep{Puebla2021.04.06.438551,kim2018not,ricci37same,messina2021solving} (including our own work) have shown that modern-day neural networks struggle to solve simple visual reasoning tasks when tested on a popular benchmark called synthetic visual reasoning test (SVRT) by \citet{fleuret2011comparing} otherwise simple for humans. We found a similar trend when we tested popular reasoning architectures like Relational Network~\citep{santoro2017simple}, Transformer~\citep{vaswani2017attention}, ESBN~\citep{Webb2021EmergentST} on Abstract Reasoning Task (ART) where the stimulus contains a simple Unicode character. As a result, visual reasoning has become an increasingly popular topic of research in recent years with the emergence of numerous fluid intelligence tests for AI algorithms, including tests for Compositional Visual Reasoning (CVR)  \citep{zerroug2022benchmark}, Ravens' (RPM) \citep{barrett2018measuring,zhang2019raven} and visual progressive matrices (V-PROM) \citep{barrett2018measuring,teney2020v} as well as an Abstract Reasoning Corpus (ARC) \citep{chollet2019measure}. 

We began this thesis by studying the computational mechanisms involved in solving the Synthetic Visual Reasoning Test (SVRT) challenge \citep{fleuret2011comparing}. This challenge consists of twenty-three binary classification tasks, each involving unique abstract relations in their formulation. Previous studies have identified two broad categories of SVRT tasks \citep{Stabinger2016,kim2018not,yihe2019program} -- tasks involving spatial-relation (\textit{SR}) judgment and tasks involving same-different (\textit{SD}) judgment. The same-different tasks are found to be harder for the neural networks compared to the spatial relation tasks \citep{Ellis2015UnsupervisedLB,kim2018not,stabinger201625,stabinger2021evaluating,Puebla2021.04.06.438551,messina2021recurrent,vaishnav2021understanding}. Consistent with this work, we proposed a novel taxonomy beyond the two primary clusters, reflecting the number of relationships used to define a particular task. A closer examination is needed to better understand the trend reflected by the neural networks in terms of accuracy and the number of relations involved in defining a particular task. An earlier study by \citet{kim2018not} has also reported that feedforward neural networks demonstrate a `straining' effect in solving tasks involving same-different relations and hypothesized that the straining effect might be because of the lack of attention. The same was also shown with a human EEG experiment by \citet{alamia2021differential} where higher activity is recorded in the lower $\beta$ band while solving the same-different judgment when compared to spatial relation judgment indicating higher demands for attention and/or working memory. To test the same, in the next chapter, we focused on understanding the role of attention in solving visual reasoning tasks.

Inspired by the two types of visual attention, we proposed a self-attention module that can be used as a \textit{feature-based} or \textit{spatial} attention to augment the features of a feedforward network (ResNet50 \citep{he2016deep}). We evaluated both types of attention-augmented neural networks on SVRT tasks and found that our proposed attentional models could solve the most challenging SVRT tasks efficiently. The relative improvements obtained by feedforward networks endowed with the two different forms of attention varied across SVRT tasks. We observed that many tasks benefited from spatial attention mechanisms, whereas a few tasks from feature-based attention and showed a significant improvement. 
Our computational analysis also leads to testable predictions for human experiments by suggesting tasks that benefit from spatial attention (task 22) or feature-based attention (task 21), tasks that benefit from either form of attention (task 19), and tasks that do not benefit from attention (task 2). While we evaluated two types of attention systems, there is a future possibility to add experiments with the third type of attention -- object-based attention \citep{duncan1984selective,egly1994shifting,vecera1994does,kramer1997object}. Object-based attention focuses on the particular object rather than its spatial location or corresponding features.


In the last part of the thesis, we proposed a novel architecture, the Guided Attention Model for (visual) Reasoning (\textit{GAMR}). We integrated both cognitive abilities humans use -- attention and memory in solving reasoning tasks. It draws inspiration from the cognitive science literature on active vision, where the spotlight of attention is routed in the visual system to gather task-relevant information. According to the theory of active vision, the visual world is explored using rapid eye movements guided by shifts of visual attention. We designed a controller akin to the mechanisms involved in the active vision framework to route the spotlight of attention and send the task-relevant representations in the memory block later used for reasoning. In GAMR, the controller is implemented with the key/query/value-based self-attention layer. Contrary to the existing method where key, query and value vectors all correspond to the same vector, the query is internally generated at each time step in our model. It helps the controller to shift the spotlight of attention. 
One of the limitations of the current approach is the fixed number of time steps. I believe that a future continuation of this work could be to incorporate a mechanism to adapt the number of time steps based on the complexity of the task. For now, we have set the number of time steps as four for all the tasks; however, a simple task might require fewer time steps to arrive at a decision with high confidence. To make the model adaptive to the situation, one possibility could be to train it with a confidence variable as a stopping criterion. 

While we have limited our analysis to synthetic visual reasoning datasets, a future possibility exists to test the models on a real-world dataset like V-PROM. It consists of images organized in a Ravens' style of reasoning with some context images and some choice images from which the correct answer is selected. Another possible direction is to think of an architecture that considers two important traits -- efficient use of data and efficient use of the computational resource. One way to design this architecture is by incorporating a read-and-write mechanism similar to a Neural Turing Machine \citep{graves2014neural}. Both these mechanisms will help the network read the already stored relations from memory and write them into the memory if they are novel. We expect such cognitive architecture to demonstrate higher-order reasoning ability, continual learning, compositionally, and meta-learnability.

We also evaluated ViT~\citep{Dosovitskiy2020-iq} -- a full self-attention architecture on SVRT tasks and found that it struggles to learn the simplest of the SVRT tasks; however, \citet{messina2021recurrent} conducted a similar study on a smaller subset of four SVRT tasks trained on 28k samples and found that a recurrent version of ViT -- an attentional network with a convolutional backbone can learn those tasks. Adding convolutions in the early layers of ViT is found to help obtain better accuracy and improve sensitivity to the optimization settings \citep{xiao2021early}. This observation motivated us to propose \textit{Conviformer} \citep{vaishnav2022conviformers} for another collaborative project on leaf-fossil classification. We propose a network to incorporate a convolutional network as the front end for a full self-attention-based vision transformer network enhancing its ability to process higher-resolution images. While bigger images hold great importance in computer vision applications like object detection, segmentation and fine-grained classification, they cannot be used with vision transformers because of the associated computational memory demand. \textit{Conviformer} improves the performance of vision transformers by incorporating local features and infusing convolutional priors in a transformer architecture. We would like to see how convolution-induced vision transformers perform on SVRT tasks. 

Concept learning is yet another exciting direction of research. One of the key features of human intelligence is the ability to quickly learn new concepts and use them to generalize to a novel scenario. A \textit{concept} can be an idea representing a class of events (e.g., walking), objects (e.g., cats), or their properties (e.g., blue color). To test the concept learning ability of neural networks in a few-shot manner, we recently introduced a novel visual reasoning dataset, Compositional Visual Reasoning (\textit{CVR})~\citep{zerroug2022benchmark}. This dataset is based on the principle of odd-one-out reasoning. In this form of reasoning task, three out of four samples follow a similar concept (rule) in their formulation, while the fourth does not. Each sample contains shapes similar to the shapes used in the SVRT challenge. It extends the variety of relations used in the formulation compared to previously defined datasets like SVRT or RPM. We have also included compositionality prior in the dataset, where some elementary relations are used to compose the several tasks. The motivation is to push the community to build a compositional and sample-efficient network.

In this thesis, we embarked on one of the pioneering endeavors to investigate self-attention through the lens of visual reasoning. Attention assumes a pivotal role in showcasing visual reasoning capabilities, and an enhanced attentional model holds promise for improved reasoning abilities. We elucidated how self-attention operations can serve as a computational model for a visual attention system, encompassing spatial and feature-based attention, while also acting as a model for active vision. Our findings indicate that self-attention is equally effective in addressing reasoning tasks as it is in tackling other challenges in the realm of vision. However, further analysis is required to unravel the underlying mechanisms within a comprehensive self-attention model, which may constrain its sample-efficient learnability for reasoning tasks. Collectively, this work exemplifies the potential advantages of incorporating self-attention mechanisms into cognitive and computer vision architectures to conquer visual reasoning tasks.